\chapter{Implementation}

This chapter maps the design decisions from Chapters 7 and 8 to concrete
repository targets and runtime responsibilities. Its purpose is traceability:
for each report claim, a reader should be able to identify where the behavior is
implemented and where supporting artifacts are produced.

The implementation is organized so benchmark reproducibility is a first-class
constraint. Each target has a stable workspace path, host binary entrypoint,
guest method binary, and benchmark ID family. This allows report statements to
be traced back to code and run artifacts without interpretation gaps.

\section{Repository Targets}

The codebase is organized into benchmark targets and crypto crates that follow a
common host/guest pattern:

\begin{itemize}
\item operation suite: \texttt{operation-bnchmrk-r0};
\item baseline ciphers: \texttt{aes-r0}, \texttt{salsa-r0}, \texttt{lowmc-r0};
\item optimization workspaces: \texttt{aes-r0-optimised},
\texttt{lowmc-r0-optimised};
\item final authenticated pipeline split:
\texttt{aes-ctr} (CTR-only) and \texttt{aes-ctr-hmac} (CTR+HMAC).
\end{itemize}

Primitive-comparison work and final authenticated-pipeline work are deliberately
kept as separate targets so experimental iterations in one area do not silently
change assumptions in another.

The CTR and CTR+HMAC split gives each mode family stable benchmark IDs for
aggregation and plot ordering, which is the implementation-specific reason for
preserving the target separation here.

Table~\ref{tab:implementation-traceability} maps the main report claims to the
code paths and artifacts that support them.

\begin{table}[H]
\centering
\small
\begin{tabular}{p{0.27\linewidth}p{0.30\linewidth}p{0.33\linewidth}}
\hline
Claim & Implementation path(s) & Evidence output(s) \\
\hline
CTR and CTR+HMAC are separate proved targets &
\path{aes-ctr/host/src/main.rs},
\path{aes-ctr-hmac/host/src/main.rs} &
campaign aggregate summaries for both target families \\
Guest binds cryptographic outputs to journal commits &
\path{aes-ctr/methods/guest/src/main.rs},
\path{aes-ctr-hmac/methods/guest/src/main.rs} &
per-target trial JSON records (raw trial outputs) \\
Proof verification is explicit in host flow &
\texttt{receipt.verify(METHOD\_ID)} call in host binaries &
stdout/stderr trial logs emitted by harness runs \\
Benchmark reproducibility is config-driven &
\path{bench-harness/config.report-main.toml},
\path{bench-harness/config.report-ops.toml} &
run manifests plus aggregate summaries per campaign \\
\hline
\end{tabular}
\caption{Implementation traceability mapping from chapter-level claims to code and artifacts.}
\label{tab:implementation-traceability}
\end{table}

\section{Guest Responsibilities}

Guest programs perform deterministic cryptographic work and produce the journal
commitments consumed by the host. Across the two final AES targets, they:

\begin{itemize}
\item deserialize keys, nonce/IV, plaintext, and AAD where applicable;
\item run AES-CTR encryption;
\item compute HMAC tag for authenticated mode;
\item commit ciphertext, tag where present, and public outputs required for
verification.
\end{itemize}

The final split keeps CTR-only and CTR+HMAC as separate guest entrypoints,
avoiding mode branches inside guest code while preserving one execution trace
per benchmark target. In authenticated mode, the guest follows the placement
rule from Chapter~8 so encryption and tag output remain part of one proved
execution.

The guest journal schema is a compatibility contract with host decode logic and
benchmark aggregation; field or serialization changes must be updated on both
sides.

\subsection*{Code Example: Guest CTR+HMAC Entrypoint}

The following excerpt shows the core guest path that performs Encrypt-then-MAC,
runs the Chapter~8 verify-before-decrypt self-check, and commits outputs to the
journal. The decrypt step here is a benchmark self-check, not a separate receiver
workflow:

\begin{verbatim}
let (ciphertext, tag) = encrypt_then_mac(
    &spec.plaintext,
    &spec.enc_key,
    &spec.mac_key,
    &spec.iv,
    &spec.aad,
);

let recovered = verify_then_decrypt(
    &ciphertext, &spec.enc_key, &spec.mac_key, &spec.iv, &spec.aad, &tag,
).expect("authentication failed in guest");

assert!(recovered == spec.plaintext);
env::commit(&AesCtrResult { ciphertext, tag });
\end{verbatim}

This snippet corresponds to
\texttt{aes-ctr-hmac/methods/guest/src/main.rs}.

\section{Host Responsibilities}

The host orchestrates proving and verification:

\begin{itemize}
\item prepare serialized guest inputs;
\item invoke prover and collect receipt;
\item verify receipt validity and image binding;
\item export benchmark-facing JSON fields for aggregation.
\end{itemize}

In authenticated mode, host-side checks follow the Chapter~8 authentication
contract. Receipt verification remains a separate host-side check that binds the
journal to the expected guest image.

For the final AES targets, host binaries also emit benchmark IDs that encode
mode family and block count (for example \texttt{aes-ctr-4blk} and
\texttt{aes-ctr-hmac-4blk}). This naming is consumed by the harness and should
remain stable across reruns to keep result comparisons valid.

The host is also responsible for benchmark-facing JSON output fields used by the
aggregation pipeline, which is why implementation and evaluation chapters are
tightly linked.

Operationally, the host is the point where experiment policy is applied:
selected block sizes, trial counts via harness config, timeout policy, and run
metadata collection. As a result, reproducibility claims depend on host-side
orchestration just as much as guest correctness.

\section{Proof Flow}

Runtime flow is:

\begin{enumerate}
\item host serializes inputs and invokes guest execution;
\item guest computes ciphertext and optional authentication tag;
\item prover generates receipt for the execution trace;
\item verifier checks receipt and journal commitments.
\end{enumerate}

This binds cryptographic outputs to a publicly verifiable execution record.
From a report perspective, it also provides the reproducibility path: campaign
metadata, target command, output JSON, and receipt verification state can be
inspected together.

The flow is identical in structure for \texttt{aes-ctr} and
\texttt{aes-ctr-hmac}, which allows direct mode-family comparisons under one
execution model. The main variable is authenticated-workload overhead, not a
different proving architecture.

\section{Reproducible Execution}

The main commands used to reproduce the implemented paths from the repository
root are:

\begin{verbatim}
make risc0-prod PROJECT=aes-ctr
make risc0-prod PROJECT=aes-ctr-hmac
python3 bench-harness/runner.py --config bench-harness/config.report-main.toml
python3 bench-harness/aggregate.py --output-root artifacts/benchmarks/report-benchmarks/main
python3 bench-harness/plot.py --output-root artifacts/benchmarks/report-benchmarks/main
\end{verbatim}

The operations companion campaign is run with:

\begin{verbatim}
python3 bench-harness/runner.py --config bench-harness/config.report-ops.toml
python3 bench-harness/aggregate.py --output-root artifacts/benchmarks/report-benchmarks/ops
\end{verbatim}

The 1--16 block CTR trace-rounding sweep used in Chapter 11 is run with:

\begin{verbatim}
python3 bench-harness/runner.py --config bench-harness/config.report-ctr-1to16.toml
python3 bench-harness/aggregate.py --output-root artifacts/benchmarks/report-benchmarks/ctr-1to16
python3 bench-harness/plot_ctr_trace_rounding.py --output-root artifacts/benchmarks/report-benchmarks/ctr-1to16
\end{verbatim}

A strict rerun should use the commit-locked report benchmark target so the run
manifest records the repository commit hash, dirty/clean state, selected config
hash, and tool versions alongside command and path provenance.

\section{Edge Cases and Failure Handling}

The implementation and harness handle failures at two levels. Cryptographic
failure, such as tag mismatch in authenticated mode, is treated as a failed
verification path and should not release unauthenticated plaintext. Benchmark
failure, such as timeout, process error, or malformed JSON, is captured by the
harness as an unsuccessful trial rather than silently entering aggregate values.

The important edge cases for the final authenticated path are:

\begin{itemize}
\item empty or very small payloads, where fixed receipt overhead dominates;
\item AAD, IV, ciphertext, or tag mismatch, where verification must fail before
decryption output is trusted;
\item block-size changes, where benchmark IDs and payload sizes must remain
consistent with harness config;
\item receipt verification failure, where the host must reject the transcript
even if ciphertext and tag have the expected byte shape.
\end{itemize}

These cases are discussed in evaluation because they affect how overhead numbers
should be interpreted. For example, HMAC cost is visible in user cycles even when
single-host wall-clock total time is noisy.
